\subsection{The literal rational specialization of scan 12}

\begin{proposition}[A-025: \status{Proved}; the all-space $T=1$ specialization]
\label{sa08:A025}
For every commutative ring $k$ and simplicial set $X$, there is a
natural isomorphism of differential graded algebras
\[
 C_{\mathrm{lit}}(X;k)/(T-1)
   \cong \Omega_{\mathrm{PL}}(X;k\otimes_{\Z}\Q),
\]
where the right side consists of compatible ordinary polynomial
differential forms on the affine simplices $\sum X_i=1$.
The assertion fails with $C_{\mathrm{gr}}$ in place of
$C_{\mathrm{lit}}$ for arbitrary infinite $X$. This verifies the
literal assertion R2 and the distinction in its attached check event
S3-136-002-141-CE008.
\end{proposition}
\begin{proof}
In each simplex and exterior degree, multiplication by $T-1$ on
$A^q$ is injective: the least nonzero weight of a finite sum survives
with coefficient $-1$. In its quotient, the identity
$n!T^{[n]}=T^n$ implies that $n!$ is invertible for every $n$.
Thus the coefficient ring becomes $k\otimes\Q$, divided powers
become ordinary powers divided by factorials, and $T=1$ imposes
exactly $\sum X_i=1$ and $\sum dX_i=0$. This identifies the
simplexwise quotient with the indicated ordinary forms, naturally in
all simplicial operators, products and the differential.

It remains to prove that taking actual sections over $X$ preserves
this quotient; right exactness of sections cannot be assumed.
As a simplicial module, $A^q$ is the sum of $\Lambda^q\Psi$ and
contractible positive-PD-degree terms. The normalized complex of
$\Lambda^q\Psi$ is $k$ in homological degree $q$ and zero
elsewhere: in the usual interval basis the exterior monomials are
degenerate except for the single monomial in simplicial degree $q$.
Hence $\pi_q(A^q)=k$, with no other homotopy. Multiplication by
$T$ sends its only surviving weight $q$ into weight $q+1$, whose
term is contractible, and therefore acts by zero on homotopy.
The exact sequence
\[
 0\longrightarrow A^q\xrightarrow{T-1}A^q
 \longrightarrow\Omega^q_{\mathrm{PL}}(-;k\otimes\Q)
 \longrightarrow0
\]
shows that the last simplicial module is weakly contractible.

The last surjection is a Kan fibration. Every map from $X$ into
its weakly contractible target is nullhomotopic. Lift a nullhomotopy
starting with the zero lift; the lifting property against
$X\times\{0\}\hookrightarrow X\times\Delta[1]$ yields an
actual section lift at the other endpoint. This proves surjectivity
on sections for every $X$, without a finiteness assumption.
The kernel is exactly $(T-1)\Gamma_X(A^q)$: the unique levelwise
preimages of a kernel section are compatible with simplicial maps.
The simplexwise algebra and differential identifications therefore
give the asserted DGA isomorphism.

For the failure with globally finite weights, take $k=\Z$ and
the discrete simplicial set $X=\mathbb N$. Its degree-zero source
quotient is
\[
 \bigl((\prod_{n\geq0}\Z)\otimes\Gamma_{\Z}(T)\bigr)/(T-1)
 \cong(\prod_{n\geq0}\Z)\otimes\Q,
\]
whose image in $\prod_{n\geq0}\Q$ consists of sequences with a
common denominator. The target contains $(1/n!)_{n\geq0}$,
which has no common denominator. Both complexes are concentrated
in degree zero for this discrete $X$, so this is a complete
counterexample.
\end{proof}
